\chapter[2017 --- Hall et al.]{2017: Jeffrey C. Hall, Michael Rosbash, Michael W. Young}
\label{ch:2017_Hall}

\section*{Main Contribution}

\textit{Discovered the molecular mechanisms controlling circadian rhythm, explaining how our internal biological clock synchronizes physiology with the 24-hour day-night cycle.}\cite{nobel-2017,lecture-2017}

\section{Official Citation}

for their discoveries of molecular mechanisms controlling circadian rhythm

\section{Background and Historical Context}

The 2017 Nobel Prize in Physiology or Medicine recognized Jeffrey C. Hall, Michael Rosbash, and Michael W. Young ``for their discoveries of molecular mechanisms controlling circadian rhythm.'' Their work revealed how the body's internal clock---the biological clock that governs the approximately 24-hour cycles of physiology and behavior known as circadian rhythms---functions at the molecular level, providing fundamental insights into one of the most ubiquitous and important features of living organisms.

Circadian rhythms, from the Latin \textit{circa diem} (``about a day''), are 24-hour biological cycles that regulate virtually every aspect of physiology and behavior in living organisms, from single-celled bacteria to humans. These rhythms control the sleep-wake cycle, hormone secretion, body temperature, metabolism, immune function, and countless other physiological processes. The existence of an internal biological clock that runs on an approximately 24-hour cycle, even in the absence of external time cues, was first demonstrated in the 1970s through experiments in which organisms were placed in constant darkness and continued to exhibit rhythmic patterns of activity and rest.

The fruit fly \textit{Drosophila melanogaster} has been a model organism for the study of circadian rhythms since the pioneering work of Konopka and Benzer in the 1970s, who identified the first clock gene mutation---\textit{period} (\textit{per})---in flies that disrupted the normal 24-hour rhythm of eclosion (emergence from the pupal case) and locomotor activity. The \textit{period} gene was subsequently cloned in 1984 by Jeffrey Hall and Michael Rosbash at Brandeis University, and its protein product (PER) was shown to oscillate with a 24-hour rhythm, rising during the night and falling during the day.

However, the molecular mechanisms by which the \textit{period} gene and its protein product generated a 24-hour rhythm remained mysterious. A fundamental question was how the molecular clock could oscillate with a period of approximately 24 hours---a timescale that is much longer than the half-lives of most proteins and mRNAs. The identification of additional clock genes and the elucidation of the feedback loops and regulatory mechanisms that generate the 24-hour oscillation became one of the major challenges in chronobiology.

\section{The Discovery/Contribution}

The elucidation of the molecular mechanisms of circadian rhythm involved the combined efforts of Hall, Rosbash, and Young, as well as many other researchers. Their contributions, made over a period of more than two decades, revealed a complex network of transcriptional feedback loops that generate and sustain the 24-hour oscillation.

\textbf{Jeffrey C. Hall and Michael Rosbash:} Hall and Rosbash, working at Brandeis University, made several key contributions to the understanding of the \textit{period} gene and its role in circadian rhythm. After cloning the \textit{period} gene in 1984, they showed that the PER protein oscillated with a 24-hour rhythm, accumulating during the night and being degraded during the day. They proposed that PER functioned within a negative feedback loop: PER protein accumulated in the cytoplasm during the night, entered the nucleus during the night, and then inhibited the transcription of its own gene, leading to a decline in PER levels during the day. When PER levels fell below a threshold, the inhibition was relieved, allowing new PER mRNA and protein to be synthesized, and the cycle began again.

This negative feedback loop model, while elegant, left a fundamental question unanswered: how could a negative feedback loop with protein synthesis, modification, nuclear import, and transcriptional repression as its components oscillate with a precise 24-hour period? Most protein half-lives are measured in hours, far shorter than 24 hours, so the feedback loop seemed too ``slow'' to generate a fast oscillation. The answer to this question required the identification of additional clock components and regulatory mechanisms that introduced time delays and buffering into the feedback loop.

Hall and Rosbash also made important contributions to the understanding of the \textit{timeless} (\textit{tim}) gene, which was identified by Hall's laboratory in 1994 through a genetic screen for clock mutants. The TIM protein was shown to interact with PER and to be required for PER's nuclear accumulation. Without TIM, PER remained in the cytoplasm and could not repress its own transcription, disrupting the negative feedback loop. The discovery of TIM revealed that PER nuclear entry was a regulated step that could serve as a point of control for the timing of the feedback loop.

A critical discovery by Hall and Rosbash was the light-dependent degradation of TIM by the photoreceptor protein CRYPTOCHROME (encoded by the \textit{cry} gene). When \textit{Drosophila} was exposed to light, CRYPTOCHROME was activated and triggered the rapid degradation of TIM, resetting the molecular clock to a new phase. This light-dependent TIM degradation provided a mechanism by which the molecular clock could be entrained (synchronized) to the external light-dark cycle, ensuring that the internal clock remained synchronized with the environmental day-night cycle. This finding explained how organisms could adjust their circadian rhythms in response to seasonal changes in day length and how jet lag occurred when the light-dark cycle was suddenly shifted.

The elucidation of the light entrainment mechanism also revealed a remarkable connection between circadian rhythms and DNA damage repair. Light exposure not only degraded TIM but also caused DNA damage, and the repair of this damage was clock-regulated. This connection suggested that circadian rhythms and DNA repair were evolutionarily linked, a finding with implications for understanding the increased cancer risk associated with shift work and circadian disruption.

\textbf{Michael W. Young:} Young, working at Rockefeller University, made several key discoveries that complemented and extended the work of Hall and Rosbash. In 1994, Young identified the \textit{doubletime} (\textit{dbt}) gene, which encoded a casein kinase 1 (CK1) enzyme that phosphorylated PER protein and promoted its degradation. Mutations in \textit{dbt} that reduced kinase activity led to the accumulation of PER and lengthened the period of the circadian clock, while mutations that increased kinase activity shortened the period. This discovery provided the first mechanistic explanation for how post-translational modifications could regulate the timing of the circadian clock.

Young also identified the \textit{clock} (\textit{Clk}) gene and its protein product (CLK), which formed a heterodimer with the CYC (Cycle) protein to activate the transcription of \textit{per} and \textit{tim}. The CLK-CYC heterodimer bound to E-box elements in the promoters of clock-controlled genes and activated their transcription during the day. PER-TIM complexes, which accumulated during the night, entered the nucleus and inhibited CLK-CYC-mediated transcription, thereby repressing their own expression and completing the negative feedback loop.

The combined work of Hall, Rosbash, and Young revealed that the circadian clock was based on an autoregulatory transcription-translation feedback loop (TTFL) consisting of positive and negative elements. The positive elements (CLK-CYC) activated the transcription of the negative elements (\textit{per} and \textit{tim}), while the negative elements (PER-TIM) inhibited the activity of the positive elements. Post-translational modifications, including phosphorylation by CK1 and other kinases, regulated the stability and nuclear localization of clock proteins, introducing the time delays necessary for 24-hour oscillation.

Subsequent research has revealed that the basic TTFL mechanism discovered in \textit{Drosophila} is conserved across evolutionarily distant organisms, including mammals, where a similar feedback loop involving the CLOCK and BMAL1 proteins (the mammalian homologs of CLK and CYC) and the PER and CRY proteins (the mammalian homologs of PER and CRYPTOCHROME) generates circadian rhythms in virtually every cell of the body.

\section{Impact on Modern Medicine}

The discoveries of Hall, Rosbash, and Young have had a significant impact on our understanding of human health and disease, and have opened up new avenues for the development of chronotherapies---treatments that take into account the timing of biological processes.

\textbf{Sleep disorders and jet lag:} The understanding of the molecular mechanisms of circadian rhythm has provided fundamental insights into the biology of sleep and the mechanisms by which the circadian clock is synchronized to the external light-dark cycle. This knowledge has informed the development of treatments for circadian rhythm sleep disorders, including shift work disorder, delayed sleep phase syndrome, and jet lag. Light therapy, melatonin supplementation, and chronotherapy (timing of medication to coincide with circadian rhythms) are now used to treat a range of circadian rhythm disorders.

\textbf{Metabolic disease:} The circadian clock plays a crucial role in regulating metabolism, and disruption of circadian rhythms has been linked to metabolic diseases including obesity, type 2 diabetes, and metabolic syndrome. Studies have shown that shift work, jet lag, and other forms of circadian disruption are associated with increased risk of metabolic disease, and that the timing of food intake is an important determinant of metabolic health. The understanding of the molecular clock has informed the development of dietary and pharmacological interventions aimed at restoring circadian rhythmicity and improving metabolic health.

\textbf{Cardiovascular disease:} Circadian rhythms influence many aspects of cardiovascular physiology, including blood pressure, heart rate, and vascular tone. The incidence of cardiovascular events, including heart attacks and strokes, exhibits a marked circadian pattern, with a peak in the early morning hours. The understanding of circadian regulation of cardiovascular function has informed the timing of cardiovascular medications and the development of chronotherapeutic strategies for preventing cardiovascular events.

\textbf{Cancer:} The circadian clock has been implicated in the regulation of cell proliferation and DNA repair, and disruption of circadian rhythms has been associated with increased cancer risk. Shift work involving circadian disruption has been classified as a probable carcinogen by the International Agency for Research on Cancer (IARC). The understanding of the molecular links between the circadian clock and cell cycle regulation has suggested new therapeutic targets for cancer and has informed the development of chronotherapeutic approaches to cancer treatment, in which the timing of chemotherapy is optimized to maximize efficacy and minimize toxicity.

\textbf{Drug development and chronopharmacology:} The discovery of the molecular clock has revealed that the expression and activity of many drug-metabolizing enzymes and drug targets oscillate with circadian rhythms. This has led to the field of chronopharmacology, which aims to optimize the timing of drug administration to maximize therapeutic efficacy and minimize side effects. The understanding of circadian regulation of drug metabolism has also informed the development of new drug delivery systems that are synchronized with the body's internal clock.

\textbf{Mental health:} Circadian rhythm disruption is a hallmark of several psychiatric disorders, including major depressive disorder, bipolar disorder, and schizophrenia. The understanding of the molecular clock has suggested new approaches to the treatment of these conditions, including the use of chronotherapy, light therapy, and pharmacological agents that target circadian clock components.

Jeffrey C. Hall, Michael Rosbash, and Michael W. Young were awarded the Nobel Prize in Physiology or Medicine in 2017 ``for their discoveries of molecular mechanisms controlling circadian rhythm.'' Their work has fundamentally changed our understanding of the biological clock and has had a lasting impact on medicine.

\section{Legacy}

The legacy of Hall, Rosbash, and Young's work extends across multiple areas of biology and medicine. Their discoveries have revealed the molecular basis of one of the most fundamental features of living organisms and have provided a framework for understanding the role of circadian rhythms in health and disease.

The identification of the core molecular clock components and the elucidation of the transcription-translation feedback loop mechanism have created a detailed and comprehensive model of how the circadian clock works. This model has been refined and extended by ongoing research, revealing additional layers of regulation, including post-translational modifications, epigenetic mechanisms, and intercellular communication.

The discovery of circadian clocks in virtually every cell of the body---a finding that was enabled by the molecular understanding of the clock developed by Hall, Rosbash, and Young---has revealed that the circadian system is far more complex and pervasive than previously appreciated. This understanding has led to the concept of ``peripheral clocks'' that operate in individual tissues and organs, coordinated by a master clock in the brain's suprachiasmatic nucleus (SCN).

The field of chronobiology continues to be one of the most active and productive areas of biomedical research, with new discoveries being made about the mechanisms of circadian rhythm, the role of the clock in health and disease, and the therapeutic potential of chronotherapy. The legacy of Hall, Rosbash, and Young's work provides the foundation for these ongoing investigations and will continue to shape the future of medicine and biology for generations to come.


\section{Modern Developments}

Hall, Rosbash, and Young's circadian rhythm work has led to modern chronobiology and sleep medicine. Understanding of clock genes has led to development of chronotherapy---timing drug administration to match circadian rhythms. Understanding of circadian disruption has revealed links to cancer, metabolic disease, and mental illness. Understanding of melatonin signaling has led to new treatments for insomnia (tasimelteon). Understanding of clock mechanisms has informed development of drugs for shift work disorder and jet lag.


\section{Bibliography}

\begin{thebibliography}{9}
\bibitem{nobel-2017}
Nobel Prize Outreach, ``The Nobel Prize in Physiology or Medicine 2017,''\emph{NobelPrize.org}.\\
\url{https://www.nobelprize.org/prizes/medicine/2017/summary/}

\bibitem{lecture-2017}
Jeffrey C. Hall, ``Nobel Lecture,''\emph{NobelPrize.org}. Winner-authored primary source; downloadable from the official Nobel Prize archive.\\
\url{https://www.nobelprize.org/prizes/medicine/2017/hall/lecture/}
\end{thebibliography}
